% Main content for QKD blueprint

This blueprint documents the formalization of \textbf{Lemma 1} from the seminal paper 
\href{https://arxiv.org/abs/quant-ph/9802025}{\emph{Unconditional Security Of Quantum Key Distribution Over Arbitrarily Long Distances}} 
by Lo and Chau (1998).

The lemma establishes a fundamental result in quantum key distribution security: 
high fidelity with a pure state implies bounded von Neumann entropy. 
This formalization is expected to attract significant interest from the quantum cryptography community, 
as it provides a machine-checked proof of a cornerstone result in the security analysis of quantum key distribution protocols.

\bigskip

\begin{center}
\fbox{\begin{minipage}{0.95\textwidth}
\vspace{0.5em}
\noindent\textbf{\large Lemma 1 from Lo \& Chau (1998)}

\vspace{0.5em}
\noindent\textit{Supplementary Note 2:}

\vspace{0.3em}
\noindent\textbf{Lemma 1} (High fidelity implies low entropy): 

If $\langle R \text{ singlets}|\rho|R \text{ singlets}\rangle > 1 - \delta$ where $\delta \ll 1$, then the von Neumann entropy 
$$S(\rho) < -(1-\delta)\log_2(1-\delta) - \delta \log_2 \frac{\delta}{(2^{2R}-1)}$$

\vspace{0.3em}
\noindent\textbf{Proof:} If $\langle R \text{ singlets}|\rho|R \text{ singlets}\rangle > 1 - \delta$, then the largest eigenvalue of the density matrix 
$\rho$ must be larger than $1-\delta$. The entropy of $\rho$ is, therefore, bounded above by that of 
$$\rho_0 = \text{diag}\left\{1-\delta, \frac{\delta}{(2^{2R}-1)}, \frac{\delta}{(2^{2R}-1)}, \ldots, \frac{\delta}{(2^{2R}-1)}\right\}$$ 

That is, $\rho_0$ is diagonal with a large entry $1-\delta$ and with the remaining probability $\delta$ equally distributed 
between the remaining $2^{2R}-1$ possibilities. \hfill Q.E.D.
\vspace{0.5em}
\end{minipage}}
\end{center}

\bigskip

\begin{definition}
\label{def:vonNeumannEntropy}
\lean{vonNeumannEntropy}
\leanok
For a Hermitian matrix $A$, the \emph{von Neumann entropy} is defined as:
$$S(A) = -\sum_i \lambda_i \log(\lambda_i)$$
where $\lambda_i$ are the eigenvalues of $A$.
\end{definition}

\begin{definition}[Majorization Ordering]
\label{def:majorization}
For two vectors $x, y \in \mathbb{R}^n$ with non-negative entries, we say $x$ is \emph{majorized} by $y$, written $x \preceq y$, if:
\begin{enumerate}
\item After sorting both vectors in decreasing order to obtain $x^{\downarrow}$ and $y^{\downarrow}$, we have:
$$\sum_{i=1}^k x^{\downarrow}_i \leq \sum_{i=1}^k y^{\downarrow}_i \quad \text{for all } k = 1, \ldots, n$$
\item The total sums are equal: $\sum_{i=1}^n x_i = \sum_{i=1}^n y_i$
\end{enumerate}
\end{definition}

\begin{definition}[Schur Convexity]
\label{def:SchurConvex}
\lean{SchurConvex}
\leanok
A function $f: \mathbb{R}^n \to \mathbb{R}$ is \emph{Schur-convex} if:
$$x \preceq y \implies f(x) \leq f(y)$$
That is, $f$ preserves the majorization ordering.
\end{definition}

\bigskip

\begin{lemma}
\label{lem:slope_comparison}
\lean{slope_comparison_tlogt}
\leanok
For $f(t) = t \cdot \log(t)$, if $a \geq c$, $b \geq d$, $a \neq b$, $c \neq d$, then:
$$\frac{f(c) - f(d)}{c - d} \leq \frac{f(a) - f(b)}{a - b}$$
\end{lemma}

\begin{lemma}
\label{lem:schur_reduction}
\lean{schur_convex_reduction_to_distinct}
\leanok
For decreasing vectors $x, y$ with $x_i \neq y_i$ for all $i$, if $x \preceq y$ then:
$$\sum_i x_i \log(x_i) \leq \sum_i y_i \log(y_i)$$
\end{lemma}

\begin{theorem}
\label{thm:schur_convex_xlogx}
\lean{schur_convex_xlogx}
\leanok
The entropy function $f(x) = \sum_i x_i \log(x_i)$ is Schur-convex.
\end{theorem}

\begin{lemma}[Eigenvalue Bound from Inner Product]
\label{lem:eigenvalue_bound}
\lean{eigenvalue_bound_eigenbasis}
\leanok
Let $\rho$ be a Hermitian matrix and $v$ a unit vector with $\|v\| = 1$. 
If $\langle v, \rho v \rangle > C$, then $\rho$ has an eigenvalue $\lambda_i > C$.
\end{lemma}

\bigskip

\begin{theorem}
\label{thm:main}
\lean{high_fidelity_implies_low_entropy_equivalent}
\leanok
\uses{def:vonNeumannEntropy, thm:schur_convex_xlogx, lem:eigenvalue_bound}
Let $\rho$ be an $R \times R$ positive semidefinite density matrix with trace 1.
If there exists a unit vector $v$ such that $\langle v, \rho v \rangle > 1 - \delta$ where $0 \leq \delta < 1$ and $\delta \leq (R-1)/R$, then:
$$S(\rho) \leq -(1-\delta)\log(1-\delta) - \delta \log\left(\frac{\delta}{R-1}\right)$$
\end{theorem}

\begin{proof}
\uses{lem:eigenvalue_bound, thm:schur_convex_xlogx}
The proof proceeds in 9 steps:

\textbf{Step 1: Eigenvalue Bound from High Fidelity.}
From the hypothesis $\langle v, \rho v \rangle > 1 - \delta$, we apply 
Lemma~\ref{lem:eigenvalue_bound} to conclude that at least one eigenvalue $\lambda_{i_0} > 1 - \delta$.

\textbf{Step 2: Eigenvalue Non-negativity.}
Since $\rho$ is positive semidefinite, all eigenvalues are non-negative: $\lambda_i \geq 0$ for all $i$.

\textbf{Step 3: Construct Comparison Distribution.}
Define a comparison eigenvalue distribution $\lambda_{\text{comp}}$:
$$\lambda_{\text{comp}}(i) = \begin{cases}
1 - \delta & \text{if } i = 0 \\
\frac{\delta}{R-1} & \text{otherwise}
\end{cases}$$
This represents a "simple" distribution with one large eigenvalue and $R-1$ equal small eigenvalues.

\textbf{Step 4: Verify Normalization.}
We verify that:
$$\sum_i \lambda_{\text{comp}}(i) = (1-\delta) + (R-1) \cdot \frac{\delta}{R-1} = (1-\delta) + \delta = 1$$

\textbf{Step 5: Establish Majorization (Key Step).}
We prove $\lambda_{\text{comp}} \preceq \lambda$ by showing:
\begin{enumerate}
\item[(a)] For all $k$: $\sum_{j \leq k} \lambda_{\text{comp}}(\sigma_j) \leq \sum_{j \leq k} \lambda(\tau_j)$ 
where $\sigma, \tau$ are sorting permutations
\item[(b)] Total sums are equal (both equal 1)
\end{enumerate}
The proof uses case analysis ($k=0$ vs $k>0$) with a proof by contradiction for the $k>0$ case.

\textbf{Step 6: Apply Schur-Convexity (Payoff).}
Using Theorem~\ref{thm:schur_convex_xlogx} and the majorization from Step 5:
$$\sum_i \lambda_i \log(\lambda_i) \geq \sum_i \lambda_{\text{comp}}(i) \log(\lambda_{\text{comp}}(i))$$

\textbf{Step 7: Unfold Entropy Definition.}
By Definition~\ref{def:vonNeumannEntropy}:
$$S(\rho) = -\sum_i \lambda_i \log(\lambda_i)$$

\textbf{Step 8: Compute Comparison Entropy.}
Direct calculation yields:
$$\sum_i \lambda_{\text{comp}}(i) \log(\lambda_{\text{comp}}(i)) = (1-\delta)\log(1-\delta) + \delta \log\left(\frac{\delta}{R-1}\right)$$

\textbf{Step 9: Final Calculation.}
Combining all steps:
\begin{align*}
S(\rho) &= -\sum_i \lambda_i \log(\lambda_i) \\
&\leq -\sum_i \lambda_{\text{comp}}(i) \log(\lambda_{\text{comp}}(i)) \\
&= -(1-\delta)\log(1-\delta) - \delta \log\left(\frac{\delta}{R-1}\right)
\end{align*}
\end{proof}

\bigskip

\noindent\textbf{Remark:} The original paper assumes $\delta \ll 1$, but for rigorous formalization we require 
$\delta \leq \frac{R-1}{R}$. This ensures $\frac{\delta}{R-1} \leq 1 - \delta$, which guarantees the comparison 
distribution $\lambda_{\text{comp}}$ is properly ordered with one dominant eigenvalue. Without this constraint, 
the majorization argument fails. Since $\frac{R-1}{R} \to 1$ as $R \to \infty$, this constraint is nearly equivalent 
to $\delta < 1$ for large systems and is always satisfied in practical QKD where $\delta \ll 1$.

